\section{Discussion}
\label{sec:discussion}

\subsection{What ``Prompt Design Matters'' Actually Means Here}
\label{sec:discussion:prompt-effects}

The statistical results (Sec.~\ref{sec:results:significance}) show a
real, measurable shift in model behavior for the large majority of this
study's prompt-design choices: phi3's role-plus-aspects-plus-few-shot
jump (+0.428 F1, V1$\to$V9) is the largest single effect this study's
significance tests measure. Fine-tuning was never put through the same
tests, but for this same model the QLoRA gain is larger still: F1 0.734
against V1's 0.138 (Table~\ref{tab:baselines}), a +0.596 improvement.
Prompt design's effect is large by the standards of training-free
methods; whether it is large by the standards of adapting the model's
own weights is a separate, smaller claim, and the baseline comparison
(Sec.~\ref{sec:results:baselines}) answers it: those effects operate
entirely within a space that, for five of six models, tops out below a
linear classifier trained on the same 14{,}000 examples in a fraction of
a second on a CPU. Prompt design is a real lever on this task, but it is
attached to a ceiling that a much simpler method already clears.

That ceiling moves, too. The few-shot-size sweep
(Sec.~\ref{sec:results:sensitivity}) shows four of six models score
higher at $k=4$ than the main grid's own $k=2$ default, so
Table~\ref{tab:main-grid} understates each model's real training-free
ceiling. Recomputing the baseline gap at each model's best available $k$
does not close it: mistral moves from $-0.050$ to $-0.022$ relative to
classical, gemma2 from $-0.087$ to $-0.076$, llama3.1 from $-0.067$ to
$-0.066$, all still below the classical baseline, while qwen3, already
ahead at $k=2$, extends its lead to $+0.028$. The direction of the
honest-limits finding survives its own most favorable adjustment; only
its margin moves.

One axis mattered less than prior in-context-learning work on other
tasks would predict. Lu et al.~\cite{lu2022fantastically} found that
reordering few-shot demonstrations alone can swing a model between
near-state-of-the-art and near-random performance. Reordering entire
prompt components produced nothing like that here: four of six models had
no meaningfully reorderable best variant at all, and where reordering did
apply, only phi3 showed a real effect ($-0.032$ F1). Component-order
stability is a genuine, if narrow, piece of good news for anyone
assembling a prompt in this space: which components to include is what
matters, not primarily what order to put them in.

\subsection{The Cost of the One Exception}
\label{sec:discussion:qwen3-cost}

qwen3 is the sole model whose best training-free prompt beats the
classical baseline on F1, and the paper's own framing needs to be honest
about what that win costs. Table~\ref{tab:cost} shows qwen3's full grid
running 5--16$\times$ longer than any other model's, driven by its native
reasoning mode left unconstrained on chain-of-thought variants
(Sec.~\ref{sec:experimental:generation}). Its margin over the classical
baseline is $+0.010$ F1; the classical baseline itself runs in a fraction
of a second per example with no GPU at all. A $+0.010$ F1 gain purchased
at 4--16$\times$ the inference cost of every other prompted model, and
still behind the classical baseline on accuracy (0.773 vs.\ 0.799), does
not read as an argument for training-free prompting over a linear
classifier in any setting where inference cost matters. Even the one
model that wins does not win by much, or cheaply.

That framing holds specifically under F1, the metric this study's main
grid and significance testing are built around; it does not hold under
every lens. Sec.~\ref{sec:results:auc}'s confidence-based AUC comparison
shows qwen3's best scorable variant losing to the classical baseline
(0.838 vs.\ 0.848), while gemma2, whose F1 (0.530) trails classical by
the widest margin of any model in this study, is the one model whose AUC
beats it. ``The one exception'' depends on which metric answers the
question, and qwen3's true best variant (V8) cannot be evaluated by the
AUC method at all. This does not weaken the cost argument above; it
means that argument is one honest answer among two.

\subsection{Safety Alignment as an Operational Risk}
\label{sec:discussion:refusal}

Llama-3.1-8B's V14 behavior (Sec.~\ref{sec:results:llama-refusal}) is
easy to misread as a parsing bug or a prompt-engineering mistake specific
to one variant. It is neither. The model reads a sexism-classification
request, once asked to reason step by step before answering, as a
request to generate hateful content rather than analyze it, and declines
on safety grounds over half the time in V14's specific configuration;
the same decline shows up, at lower but still substantial rates, across
every other chain-of-thought variant this model was evaluated on
(Table~\ref{tab:llama-cot}), so chain-of-thought reasoning itself, not
any one combination of surrounding components, is what triggers it.
Scored under this study's own default convention, that decline registers
as an unremarkable-looking F1 of 0.330 for V14, and quietly deflated
numbers for several other cells that stay in the main table, each
indistinguishable on its face from an ordinary weak variant. Only
inspecting the raw generations, rather than trusting the aggregate
metric, revealed what was actually happening; a related but mechanically
distinct pattern (Sec.~\ref{sec:results:llama-refusal}) shows that not
every elevated \texttt{fail\_ratio} in this model is a refusal at all,
some are a fixed reasoning-token budget proving too short once the
prompt also asks for aspect-by-aspect analysis, which is a generation
setting worth revisiting rather than a safety-alignment story. Röttger et
al.'s XSTest~\cite{rottger2023xstest} showed the exaggerated-safety
version of this failure mode under adversarial probing; llama3.1's
refusals show it can also emerge unprompted, from an ordinary
classification pipeline, on genuinely on-topic content. For any
deployment that scores a refusal as a default label rather than flagging
it, this is a silent failure mode: the system does not error, it
produces a plausible-looking, systematically wrong answer on a
substantial share of a specific input class.
\texttt{fail\_ratio} (Sec.~\ref{sec:experimental:metrics}), not accuracy
or F1 alone, is what surfaces it, and only because this study checked.

\subsection{What a Single Aggregate Score Hides}
\label{sec:discussion:aggregate}

The subtype breakdown and the significance-testing disagreement point at
the same underlying problem. The subtype breakdown
(Sec.~\ref{sec:results:subtype}) shows a model's own best variant can
still miss over half of the rarest sexism category, prejudiced
discussions, while its aggregate F1 looks unremarkable: gemma2 and
phi3 post worst-variant false-negative spreads of 16\% and 96--100\%
respectively, a gap the main grid's F1 column alone never reveals. The
significance-testing disagreement (Sec.~\ref{sec:results:significance})
reflects the same class-imbalance sensitivity: phi3's largest F1 effect
in the entire study is invisible to accuracy-based McNemar testing,
because a 76\%-negative test set can absorb a large shift in
positive-class detection without moving overall accuracy much at all.
Röttger et al.'s HateCheck~\cite{rottger2020hatecheck} made the general
version of this argument for hate speech detection: held-out accuracy
and F1 can look strong while hiding specific, systematic weaknesses that
only a finer-grained test reveals. This study's subtype table and its
dual significance tests apply that same argument to sexism detection, in
two independent places rather than one.

Together, these four points sharpen this study's central claim: prompt
design produces a genuine, well-measured effect on this task, but one
that operates inside a ceiling a simple non-LLM baseline already
exceeds for most models tested; the one model that clears that ceiling
pays a steep,
possibly disqualifying compute cost to do so; and the aggregate numbers
behind all of these conclusions can themselves hide model-specific
failure modes, a safety refusal or a rare-subtype blind spot, that only
closer inspection surfaces.
